\documentclass[11pt,a4paper]{article}
% --- Pacotes Básicos e Suporte Matemático ---
\usepackage[utf8]{inputenc}
\usepackage[portuguese]{babel}
\usepackage{amsmath}
\usepackage{amsfonts}
\usepackage{amssymb}
\usepackage{graphicx}
\usepackage{booktabs}
\usepackage{geometry}
\usepackage{hyperref}

% --- Suporte a Blocos de Código ---
\usepackage{listings}
\usepackage{xcolor}
\usepackage{cite}

% --- Configuração da Página ---
\geometry{a4paper, total={170mm,257mm}, left=20mm, top=20mm}

\definecolor{codegreen}{rgb}{0,0.6,0}
\definecolor{codegray}{rgb}{0.5,0.5,0.5}
\definecolor{codepurple}{rgb}{0.58,0,0.82}
\definecolor{backcolour}{rgb}{0.95,0.95,0.92}

% --- Configuração de Estilo para Blocos de Código ---
\lstdefinestyle{mystyle}{
    backgroundcolor=\color{backcolour},   
    commentstyle=\color{codegreen},
    keywordstyle=\color{magenta},
    numberstyle=\tiny\color{codegray},
    stringstyle=\color{codepurple},
    basicstyle=\ttfamily\footnotesize,
    breakatwhitespace=false,         
    breaklines=true,                 
    captionpos=b,                    
    keepspaces=true,                 
    numbers=left,                    
    numbersep=5pt,                  
    showspaces=false,                
    showstringspaces=false,
    showtabs=false,                  
    tabsize=2
}
\lstset{style=mystyle}

% --- Informações do Documento ---
\title{\textbf{Graph-Based Cognitive Memory: Otimizando o Uso de Tokens em Agentes de IA via Neo4j, Redis e Mecanismos de Atenção Exógenos}}
\author{\textbf{Devs Café Community} \\ \small Inteligência Artificial \& Engenharia de Software}
\date{2026}

\begin{document}

\maketitle

\begin{abstract}
Modelos de Linguagem de Grande Porte (LLMs) enfrentam limitações severas de custo e processamento devido ao crescimento quadrático da computação associado à janela de contexto. Este trabalho propõe uma arquitetura de memória híbrida bimodal (Curto e Longo Prazo) que substitui a busca semântica tradicional baseada em arquivos planos por um ecossistema integrado de \textbf{Grafos de Conhecimento (Neo4j)} e \textbf{Caches de Alta Velocidade (Redis)}. Adicionalmente, introduz-se um formalismo matemático para um \textbf{Mecanismo de Atenção Exógeno Directionado à Tarefa}, operando diretamente na camada de recuperação de dados. Demonstramos que a habilidade nativa dos LLMs em gerar consultas estruturadas (Cypher) combinada a filtros de atenção matemáticos reduz drasticamente o consumo de tokens, elimina o fenômeno \textit{Lost in the Middle} e minimiza a latência de inferência.
\end{abstract}

\small{\textbf{Palavras-chave:} Grafos de Conhecimento, Neo4j, Redis, LLM, Agentes de IA, RAG, Atenção Exógena, Gestão de Tokens.}

\tableofcontents
\newpage

---

\section{Introdução}
A arquitetura tradicional de \textit{Retrieval-Augmented Generation} (RAG) baseada em bancos de dados vetoriais puros enfrenta dois gargalos fundamentais no desenvolvimento de agentes autônomos:
\begin{enumerate}
    \item \textbf{Ineficiência de Tokens:} A injeção de documentos inteiros ou múltiplos \textit{chunks} densos de texto satura a janela de contexto com redundâncias e ruídos textuais.
    \item \textbf{Perda de Contexto Relacional Global:} Fragmentos isolados de texto perdem as conexões causais, estruturais e de rede que ligam diferentes entidades informacionais ao longo do tempo.
\end{enumerate}

Ao transicionar para um modelo de \textbf{Memória Baseada em Grafos}, os dados são estruturados como entidades (Nós) e interações (Arestas), permitindo uma recuperação cirúrgica do conhecimento. Este artigo detalha os fundamentos operacionais, os modelos matemáticos de controle e os algoritmos de filtragem por atenção exógena necessários para potencializar sistemas multi-agentes.

---

\section{Componentes da Arquitetura de Memória Bimodal}

\subsection{Memória de Longo Prazo: Neo4j vs. Arquivos Planos}
Buscar informações cruzando diretórios de arquivos textuais ou partições vetoriais massivas gera alta complexidade computacional e imprevisibilidade no consumo de tokens. O Neo4j supera essa barreira por meio da indexação topológica por proximidade. 

Quando um agente interage com um conceito, o grafo expõe de forma determinística as propriedades locais adjacentes (ex: prazos, dependências, restrições). Além disso, o grafo opera a consolidação de entidades: interações repetidas sobre o mesmo tema convergem para os mesmos nós centrais, eliminando a duplicação de dados textuais no prompt do modelo.

\subsection{Memória de Curto Prazo: Redis como Espaço de Rascunho Cognitivo}
Informações voláteis e de curtíssima validade (como códigos de verificação, estados intermediários de execução ou instruções momentâneas do usuário) poluem o histórico da conversa se forem mantidas diretamente no prompt.

O Redis atua como a memória RAM de acesso rápido do sistema operacional do agente. Através de mecanismos de \textit{Function Calling}, a própria IA determina a gravação de dados temporários sob políticas de expiração estritas (\textit{Time-to-Live} - TTL), aliviando a carga contextual interna.

---

\section{Formalismo Matemático para Controle de Memória}

Para evitar a saturação de ambos os bancos de dados e garantir a eficiência econômica do sistema, a transição entre as camadas de memória é governada por funções matemáticas estritas.

\subsection{Algoritmo de Decisão de Roteamento (Curto vs. Longo Prazo)}
Toda nova unidade de informação $x$ gerada em um tempo $t$ é submetida a uma \textbf{Função de Utilidade da Informação ($U$)}:

\begin{equation}
U(x, t) = \alpha \cdot F(x) + \beta \cdot C(x) - \gamma \cdot \Delta t
\end{equation}

Onde:
\begin{itemize}
    \item $F(x) \in [0, 1]$ é a frequência estimada ou real de acessos ao conceito $x$.
    \item $C(x) \in [0, 1]$ é a métrica de centralidade de conexão do nó no grafo existente.
    \item $\Delta t = t - t_{\text{criacao}}$ é a distância temporal desde o registro do evento.
    \item $\alpha, \beta, \gamma$ são hiperparâmetros de ponderação do sistema.
\end{itemize}

A política de roteamento executa as seguintes condições de contorno:
\begin{equation}
\text{Destino}(x) = 
\begin{cases} 
\text{Persistir no Neo4j}, & \text{se } U(x, t) > \tau_{\text{high}} \\
\text{Manter no Redis Cache}, & \text{se } \tau_{\text{low}} \le U(x, t) \le \tau_{\text{high}} \\
\text{Descarte Contextual (Esquecimento)}, & \text{se } U(x, t) < \tau_{\text{low}}
\end{cases}
\end{equation}

\subsection{Otimização de Cache via LFU com Decaimento Temporal (Aging)}
Para mitigar a ineficiência de políticas LRU convencionais frente a padrões de picos estocásticos de diálogos, adota-se o LFU com \textit{Aging}. A frequência ajustada $F_{\text{adj}}$ é calculada via média móvel exponencial:

\begin{equation}
F_{\text{adj}}(t) = F_{\text{adj}}(t-1) \cdot e^{-\lambda \Delta t} + 1
\end{equation}

Onde $\lambda$ representa a taxa de decaimento do esquecimento e $\Delta t$ o intervalo de tempo desde a última requisição. O TTL dinâmico no Redis é parametrizado por:

\begin{equation}
\text{TTL}(x) = \text{TTL}_{\text{base}} \cdot \ln(1 + F_{\text{adj}}(t))
\end{equation}

\subsection{Recuperação por Subgrafos Esparsos sob Restrição de Orçamento}
Para impedir o fenômeno de \textit{Graph Explosion} (onde a busca por caminhos adjacentes consome tokens em excesso), a busca restringe a extração a um subgrafo $S$ cuja soma total de tokens das propriedades não exceda um orçamento linear rígido $\Omega$:

\begin{equation}
\max \sum_{v \in S} \text{Relevância}(v) \quad \text{sujeito a} \quad \sum_{v \in S} \text{Tokens}(v) \le \Omega
\end{equation}

---

\section{Mecanismo de Atenção Exógena Direcionada à Tarefa}

Injetar grandes volumes de dados no prompt causa a degradação da atenção intrínseca do LLM. Propõe-se um \textbf{Mecanismo de Atenção Exógeno} executado fora da rede neural do modelo de linguagem, reduzindo a complexidade de atenção quadrática original para filtros lineares na camada de dados.

O peso de atenção exógena $A$ de um nó $v$ em relação a uma tarefa específica do agente $T$ e ao histórico recente do chat $H$ é definido através do Scaled Dot-Product Attention adaptado:

\begin{equation}
A(v, T, H) = \text{Softmax} \left( \frac{Q(T, H) \cdot K(v)^T}{\sqrt{d}} \right) \cdot W_c(v)
\end{equation}

Onde:
\begin{itemize}
    \item $Q(T, H)$ (\textbf{Query}): Vetor de embedding gerado a partir da descrição da tarefa e contexto operacional recente do agente.
    \item $K(v)$ (\textbf{Key}): Vetor de embedding contendo as propriedades intrínsecas e metadados do nó $v$ (ou chave de cache).
    \item $\sqrt{d}$: Fator de escala baseado na dimensionalidade vetorial $d$ para estabilização do gradiente.
    \item $W_c(v)$ (\textbf{Peso de Contexto Causal}): Multiplicador topológico proporcional à Centralidade de Grau $C_D(v) = \deg(v)$ no Neo4j.
\end{itemize}

O operador Softmax garante que $\sum_v A(v, T, H) = 1$. Um limiar estrito de corte $\theta$ é estabelecido, de forma que o nó $v$ é injetado no prompt final apenas se $A(v, T, H) > \theta$.

---

\section{Implementação Algorítmica}

O seguinte algoritmo demonstra o fluxo de condensação contextual orientado por atenção e restrição orçamentária:

\begin{lstlisting}[language=Python, caption=Middleware de Gerenciamento Cognitivo de Contexto]
import numpy as np

def calcular_softmax(valores):
    exp_v = np.exp(valores - np.max(valores))
    return exp_v / exp_v.sum()

def condensar_contexto_por_atencao(task_desc, hist_curto, nos_neo4j, limite_tokens=600):
    # 1. Gerar o vetor Query (Q) baseado na Task do Bot
    vetor_query = obter_embedding(f"Task: {task_desc} | Historico: {hist_curto}")
    
    candidatos_memoria = []
    chaves_e_valores = nos_neo4j + buscar_chaves_recentes_redis()
    
    scores_brutos = []
    dimensao_d = len(vetor_query)
    
    # 2. Calcular o score bruto via Produto Escalar Escalonado
    for item in chaves_e_valores:
        vetor_key = obter_embedding(item.texto_para_indexacao)
        score = np.dot(vetor_query, vetor_key) / np.sqrt(dimensao_d)
        
        # Ponderacao topologica exogena (Neo4j)
        if item.origem == "Neo4j":
            score *= (1 + item.centralidade_grau * 0.1)
            
        scores_brutos.append(score)
    
    # 3. Aplicar Softmax para distribuicao de Atencao
    pesos_atencao = calcular_softmax(scores_brutos)
    
    for i, item in enumerate(chaves_e_valores):
        item.atencao = pesos_atencao[i]
        candidatos_memoria.append(item)
        
    candidatos_memoria.sort(key=lambda x: x.atencao, reverse=True)
    
    # 4. Preenchimento dinamico do orcamento de tokens (\Omega)
    contexto_otimizado = []
    tokens_acumulados = 0
    
    for item in candidatos_memoria:
        if tokens_acumulados + item.tokens > limite_tokens:
            break
        contexto_otimizado.append(item.conteudo_real)
        tokens_acumulados += item.tokens
        
    return contexto_otimizado
\end{lstlisting}

---

\section{Conclusão e Trabalhos Futuros}
A integração de Neo4j e Redis estruturada sob filtros de atenção exógenos matematizados altera o paradigma de desenvolvimento de sistemas baseados em LLM. Ao delegar o processamento de busca semântica relacional e o gerenciamento de estado para algoritmos dedicados fora do transformador, o agente de IA passa a agir como uma unidade de processamento pura e focada. Como trabalhos futuros, sugere-se o estudo empírico da calibração dinâmica dos limiares de corte $\theta$ sob variações abruptas de tarefas complexas.

\end{document}
